\documentclass[12pt,a4paper]{article}

\usepackage[T5]{fontenc}
\usepackage[utf8]{inputenc}
\usepackage[vietnamese]{babel}
\usepackage{amsmath, amssymb, amsfonts}
\usepackage{booktabs}
\usepackage[margin=2.2cm]{geometry}
\usepackage{tikz}
\usepackage{xcolor}
\usetikzlibrary{positioning, arrows.meta}

\setlength{\parindent}{0pt}
\setlength{\parskip}{4pt}
\allowdisplaybreaks

% ---- Ký hiệu tắt ----
\newcommand{\relu}{\operatorname{ReLU}}
\newcommand{\softmax}{\operatorname{softmax}}
\newcommand{\pd}[2]{\dfrac{\partial #1}{\partial #2}}

% ---- Sơ đồ MLP: 3 đầu vào, 4 neuron ẩn ReLU, tầng ra tùy bài ----
% #1: danh sách "chỉ số/toạ độ y" của neuron đầu ra
% #2: nhãn hàm kích hoạt đầu ra
\newcommand{\mlpdiagram}[2]{%
\begin{tikzpicture}[x=3.4cm, y=1.0cm, >=Stealth,
  neuron/.style={circle, draw, minimum size=0.9cm, inner sep=0pt, font=\small},
  inp/.style={neuron, fill=blue!10},
  hid/.style={neuron, fill=orange!20},
  outn/.style={neuron, fill=green!20}]
  % tầng vào
  \foreach \i/\yy in {1/1, 2/0, 3/-1} {
    \node[inp] (in-\i) at (0,\yy) {$x_{\i}$};
  }
  % tầng ẩn
  \foreach \j/\yy in {1/1.5, 2/0.5, 3/-0.5, 4/-1.5} {
    \node[hid] (h-\j) at (1,\yy) {$a^{(1)}_{\j}$};
  }
  % tầng ra
  \foreach \k/\yy in {#1} {
    \node[outn] (o-\k) at (2,\yy) {$\hat y_{\k}$};
  }
  % kết nối
  \foreach \i in {1,2,3} {
    \foreach \j in {1,2,3,4} {
      \draw[->, gray!70] (in-\i) -- (h-\j);
    }
  }
  \foreach \j in {1,2,3,4} {
    \foreach \k/\yy in {#1} {
      \draw[->, gray!70] (h-\j) -- (o-\k);
    }
  }
  % nhãn tầng
  \node[font=\footnotesize, align=center] at (0,-3.1) {Đầu vào\\3 neuron};
  \node[font=\footnotesize, align=center] at (1,-3.1) {Tầng ẩn\\4 neuron\\ReLU};
  \node[font=\footnotesize, align=center] at (2,-3.1) {Đầu ra\\#2};
\end{tikzpicture}%
}

\begin{document}

\begin{center}
  {\Large\bfseries LỜI GIẢI BÀI TẬP MLP}\\[2pt]
  {\large Học sâu --- Thiết kế MLP và tính toán thủ công}\\[2pt]
  \textit{Mock MLP 2026-10-07 \quad (90 phút, 10 điểm)}
\end{center}

\section*{Quy ước chung}
\begin{itemize}
  \item Mọi vector là vector cột. $W^{(l)}$ có số hàng bằng số neuron của tầng hiện tại và số cột bằng số đầu vào của tầng đó. Mỗi neuron ở tầng ẩn và tầng ra có một bias riêng.
  \item Truyền xuôi: $z^{(l)} = W^{(l)} a^{(l-1)} + b^{(l)}$, $a^{(l)} = f(z^{(l)})$, với $a^{(0)} = x$.
  \item $\ln$ là logarit tự nhiên. Kết quả không hữu tỉ giữ $\ge 6$ chữ số thập phân ở bước trung gian, kết quả cuối làm tròn đến 4 chữ số thập phân.
  \item $\relu(z)=\max(0,z)$, $\relu'(z)=1$ nếu $z>0$ và $\relu'(z)=0$ nếu $z\le 0$ (quy ước $\relu'(0)=0$).
\end{itemize}

% =====================================================================
\section*{Câu 1. Thiết kế kiến trúc MLP (2,00 điểm)}

Mỗi mẫu có $n_0 = 3$ đặc trưng thực; MLP kết nối đầy đủ với một tầng ẩn gồm $n_1=4$ neuron ReLU.

\subsection*{a) Bài toán A: dự đoán nhiệt độ (hồi quy, giá trị âm hoặc dương)}

\begin{center}
\begin{tabular}{@{}ll@{}}
\toprule
Thành phần & Lựa chọn \\
\midrule
Số chiều đầu vào & $3$ \\
Số neuron ẩn / kích hoạt ẩn & $4$ / ReLU \\
Số neuron đầu ra & $1$ \\
Kích hoạt đầu ra & \textbf{Tuyến tính (identity)}: $\hat y = z^{(2)}$ \\
Hàm mất mát & \textbf{MSE}: $L=\tfrac12(\hat y-y)^2$ \\
\bottomrule
\end{tabular}
\end{center}

\textit{Lý do:} nhiệt độ nhận giá trị thực tùy ý (âm hoặc dương) nên đầu ra không được bị chặn. Sigmoid chặn trong $(0,1)$ và ReLU không cho giá trị âm, cả hai đều không phù hợp. MSE là hàm mất mát chuẩn cho hồi quy.

\begin{center}
\mlpdiagram{1/0}{1 neuron, tuyến tính}
\end{center}

\subsection*{b) Bài toán B: phân loại thư spam ($y=1$) / không spam ($y=0$)}

\begin{center}
\begin{tabular}{@{}ll@{}}
\toprule
Thành phần & Lựa chọn \\
\midrule
Số chiều đầu vào & $3$ \\
Số neuron ẩn / kích hoạt ẩn & $4$ / ReLU \\
Số neuron đầu ra & $1$ \\
Kích hoạt đầu ra & \textbf{Sigmoid}: $\hat y=\sigma(z^{(2)})=\dfrac{1}{1+e^{-z^{(2)}}}\in(0,1)$ \\
Hàm mất mát & \textbf{Binary cross-entropy}: $L=-\big[y\ln\hat y+(1-y)\ln(1-\hat y)\big]$ \\
\bottomrule
\end{tabular}
\end{center}

\textit{Lý do:} phân loại nhị phân cần một xác suất $P(y=1\mid x)\in(0,1)$, nên dùng sigmoid. Hàm mất mát tương ứng là binary cross-entropy (cực đại hợp lý Bernoulli).

\begin{center}
\mlpdiagram{1/0}{1 neuron, sigmoid}
\end{center}

\subsection*{c) Bài toán C: phân loại vào đúng một trong ba lớp loại trừ nhau}

\begin{center}
\begin{tabular}{@{}ll@{}}
\toprule
Thành phần & Lựa chọn \\
\midrule
Số chiều đầu vào & $3$ \\
Số neuron ẩn / kích hoạt ẩn & $4$ / ReLU \\
Số neuron đầu ra & $3$ (mỗi lớp một neuron) \\
Kích hoạt đầu ra & \textbf{Softmax}: $\hat y_k=\dfrac{e^{z^{(2)}_k}}{\sum_{j=1}^{3}e^{z^{(2)}_j}}$, $\sum_k\hat y_k=1$ \\
Hàm mất mát & \textbf{Cross-entropy}: $L=-\sum_{k=1}^{3}y_k\ln\hat y_k$ \\
\bottomrule
\end{tabular}
\end{center}

\textit{Lý do:} các lớp loại trừ nhau nên đầu ra phải là một phân phối xác suất trên 3 lớp (tổng bằng 1), tức softmax. Cross-entropy là hàm mất mát tương ứng với nhãn one-hot.

\begin{center}
\mlpdiagram{1/1, 2/0, 3/-1}{3 neuron, softmax}
\end{center}

\subsection*{d) Số tham số của mô hình bài toán C}

Với $n_0=3$, $n_1=4$, $n_2=3$:

\begin{center}
\begin{tabular}{@{}lccc@{}}
\toprule
Tham số & Kích thước & Quy tắc & Số phần tử \\
\midrule
$W^{(1)}$ & $4\times 3$ & $n_1\times n_0$ & $12$ \\
$b^{(1)}$ & $4\times 1$ & $n_1\times 1$ & $4$ \\
$W^{(2)}$ & $3\times 4$ & $n_2\times n_1$ & $12$ \\
$b^{(2)}$ & $3\times 1$ & $n_2\times 1$ & $3$ \\
\bottomrule
\end{tabular}
\end{center}

\[
N_{\text{tham số}} = (4\cdot 3 + 4) + (3\cdot 4 + 3) = 16 + 15 = \boxed{31}
\]
(trong đó $24$ trọng số và $7$ bias).

% =====================================================================
\section*{Câu 2. MLP hồi quy: một bước huấn luyện (5,00 điểm)}

\textbf{Dữ liệu:}
$x=\begin{bmatrix}1\\2\end{bmatrix}$, $y=2$, $\eta=0.1$,
\[
W^{(1)}=\begin{bmatrix}1&0\\-1&1\end{bmatrix},\quad
b^{(1)}=\begin{bmatrix}0\\-2\end{bmatrix},\quad
W^{(2)}=\begin{bmatrix}0.5&-1\end{bmatrix},\quad
b^{(2)}=0.5 .
\]
$\hat y=z^{(2)}$ và $L=\tfrac12(\hat y-y)^2$.

\subsection*{a) Truyền xuôi (1,00 điểm)}

\textbf{Tầng ẩn.} Công thức: $z^{(1)}=W^{(1)}x+b^{(1)}$.
\[
z^{(1)}=\begin{bmatrix}1\cdot1+0\cdot2\\(-1)\cdot1+1\cdot2\end{bmatrix}+\begin{bmatrix}0\\-2\end{bmatrix}
=\begin{bmatrix}1\\1\end{bmatrix}+\begin{bmatrix}0\\-2\end{bmatrix}
=\begin{bmatrix}1\\-1\end{bmatrix}.
\]
\[
a^{(1)}=\relu(z^{(1)})=\begin{bmatrix}\max(0,1)\\\max(0,-1)\end{bmatrix}=\begin{bmatrix}1\\0\end{bmatrix}.
\]

\textbf{Tầng ra.} $z^{(2)}=W^{(2)}a^{(1)}+b^{(2)}$:
\[
z^{(2)}=0.5\cdot1+(-1)\cdot0+0.5=1,\qquad \hat y=z^{(2)}=1.
\]

\textbf{Mất mát.}
\[
L=\tfrac12(\hat y-y)^2=\tfrac12(1-2)^2=\tfrac12\cdot1=\boxed{0.5000}.
\]

\subsection*{b) Gradient tầng ra (0,75 điểm)}

Công thức: $\delta^{(2)}=\pd{L}{z^{(2)}}=\pd{L}{\hat y}\cdot\pd{\hat y}{z^{(2)}}=(\hat y-y)\cdot1$.
\[
\delta^{(2)}=1-2=\boxed{-1}.
\]
Vì $z^{(2)}=W^{(2)}a^{(1)}+b^{(2)}$ nên $\pd{L}{W^{(2)}}=\delta^{(2)}\,(a^{(1)})^{\top}$ và $\pd{L}{b^{(2)}}=\delta^{(2)}$:
\[
\pd{L}{W^{(2)}}=(-1)\begin{bmatrix}1&0\end{bmatrix}=\boxed{\begin{bmatrix}-1&0\end{bmatrix}}\ (1\times2),\qquad
\pd{L}{b^{(2)}}=\boxed{-1}.
\]

\subsection*{c) Gradient tầng ẩn (1,50 điểm)}

Công thức lan truyền ngược:
\[
\delta^{(1)}=\pd{L}{z^{(1)}}=\big((W^{(2)})^{\top}\delta^{(2)}\big)\odot\relu'(z^{(1)}).
\]
Tính từng thành phần:
\[
(W^{(2)})^{\top}\delta^{(2)}=\begin{bmatrix}0.5\\-1\end{bmatrix}(-1)=\begin{bmatrix}-0.5\\1\end{bmatrix},\qquad
\relu'(z^{(1)})=\begin{bmatrix}\relu'(1)\\\relu'(-1)\end{bmatrix}=\begin{bmatrix}1\\0\end{bmatrix},
\]
\[
\delta^{(1)}=\begin{bmatrix}-0.5\\1\end{bmatrix}\odot\begin{bmatrix}1\\0\end{bmatrix}=\boxed{\begin{bmatrix}-0.5\\0\end{bmatrix}}.
\]
Với $z^{(1)}=W^{(1)}x+b^{(1)}$: $\pd{L}{W^{(1)}}=\delta^{(1)}x^{\top}$, $\pd{L}{b^{(1)}}=\delta^{(1)}$.
\[
\pd{L}{W^{(1)}}=\begin{bmatrix}-0.5\\0\end{bmatrix}\begin{bmatrix}1&2\end{bmatrix}
=\boxed{\begin{bmatrix}-0.5&-1\\0&0\end{bmatrix}},\qquad
\pd{L}{b^{(1)}}=\boxed{\begin{bmatrix}-0.5\\0\end{bmatrix}}.
\]

\textbf{Quy tắc chuỗi cho $w^{(1)}_{11}$.} $w^{(1)}_{11}$ ảnh hưởng $z^{(1)}_1$, rồi $a^{(1)}_1$, rồi $z^{(2)}$, rồi $\hat y$, rồi $L$:
\[
\pd{L}{w^{(1)}_{11}}
=\pd{L}{\hat y}\cdot\pd{\hat y}{z^{(2)}}\cdot\pd{z^{(2)}}{a^{(1)}_1}\cdot\pd{a^{(1)}_1}{z^{(1)}_1}\cdot\pd{z^{(1)}_1}{w^{(1)}_{11}}.
\]
Các đạo hàm cục bộ:
\begin{align*}
\pd{L}{\hat y}&=\hat y-y=-1, &
\pd{\hat y}{z^{(2)}}&=1, &
\pd{z^{(2)}}{a^{(1)}_1}&=w^{(2)}_{1}=0.5,\\
\pd{a^{(1)}_1}{z^{(1)}_1}&=\relu'(1)=1, &
\pd{z^{(1)}_1}{w^{(1)}_{11}}&=x_1=1.
\end{align*}
\[
\pd{L}{w^{(1)}_{11}}=(-1)\cdot1\cdot0.5\cdot1\cdot1=\boxed{-0.5},
\]
khớp với phần tử hàng 1, cột 1 của $\pd{L}{W^{(1)}}$ ở trên.

\textbf{Neuron ẩn ``chết''.} $z^{(1)}_2=-1<0$ nên $\relu'(z^{(1)}_2)=0$, suy ra $\delta^{(1)}_2=0$. \textbf{Neuron ẩn thứ 2} không truyền gradient trong lượt này, và hàng 2 của $\pd{L}{W^{(1)}}$ và $\pd{L}{b^{(1)}}$ bằng $0$. Neuron ẩn 1 có $z^{(1)}_1=1>0$ nên vẫn hoạt động.

\subsection*{d) Cập nhật SGD, $\eta=0.1$ (1,00 điểm)}

Tất cả gradient ở trên được tính từ bộ tham số ban đầu. Công thức: $\theta_{\text{new}}=\theta-\eta\,\pd{L}{\theta}$.

\begin{align*}
W^{(1)}_{\text{new}}&=\begin{bmatrix}1&0\\-1&1\end{bmatrix}-0.1\begin{bmatrix}-0.5&-1\\0&0\end{bmatrix}
=\begin{bmatrix}1.05&0.1\\-1&1\end{bmatrix},\\[2pt]
b^{(1)}_{\text{new}}&=\begin{bmatrix}0\\-2\end{bmatrix}-0.1\begin{bmatrix}-0.5\\0\end{bmatrix}
=\begin{bmatrix}0.05\\-2\end{bmatrix},\\[2pt]
W^{(2)}_{\text{new}}&=\begin{bmatrix}0.5&-1\end{bmatrix}-0.1\begin{bmatrix}-1&0\end{bmatrix}
=\begin{bmatrix}0.6&-1\end{bmatrix},\\[2pt]
b^{(2)}_{\text{new}}&=0.5-0.1(-1)=0.6.
\end{align*}

\subsection*{e) Đánh giá sau cập nhật (0,75 điểm)}

\[
z^{(1)}_{\text{new}}=W^{(1)}_{\text{new}}x+b^{(1)}_{\text{new}}
=\begin{bmatrix}1.05\cdot1+0.1\cdot2\\-1\cdot1+1\cdot2\end{bmatrix}+\begin{bmatrix}0.05\\-2\end{bmatrix}
=\begin{bmatrix}1.25\\1\end{bmatrix}+\begin{bmatrix}0.05\\-2\end{bmatrix}
=\begin{bmatrix}1.3\\-1\end{bmatrix},
\]
\[
a^{(1)}_{\text{new}}=\relu(z^{(1)}_{\text{new}})=\begin{bmatrix}1.3\\0\end{bmatrix},
\]
\[
z^{(2)}_{\text{new}}=0.6\cdot1.3+(-1)\cdot0+0.6=0.78+0.6=1.38,\qquad
\hat y_{\text{new}}=1.38,
\]
\[
L_{\text{new}}=\tfrac12(1.38-2)^2=\tfrac12(0.62)^2=\tfrac12\cdot0.3844=\boxed{0.1922}.
\]
\[
\Delta L=L_{\text{new}}-L=0.1922-0.5=\boxed{-0.3078}.
\]
\textbf{Kết luận:} $\Delta L<0$, nên loss của mẫu \textbf{giảm} sau bước cập nhật (từ $0.5000$ xuống $0.1922$), và $\hat y$ tiến gần hơn đến $y=2$ (từ $1$ lên $1.38$).

% =====================================================================
\section*{Câu 3. MLP phân loại: Softmax và cross-entropy (3,00 điểm)}

\textbf{Dữ liệu:} $x=\begin{bmatrix}1\\-1\end{bmatrix}$, $y=\begin{bmatrix}0\\1\\0\end{bmatrix}$ (lớp 2),
\[
W^{(1)}=\begin{bmatrix}1&0\\0&-1\end{bmatrix},\ b^{(1)}=\begin{bmatrix}0\\0\end{bmatrix},\quad
W^{(2)}=\begin{bmatrix}1&0\\0&0\\0&1\end{bmatrix},\ b^{(2)}=\begin{bmatrix}0\\0\\0\end{bmatrix}.
\]
$L=-\sum_{k=1}^{3}y_k\ln\hat y_k$.

\subsection*{a) $z^{(1)}$, $a^{(1)}$ và logits $z^{(2)}$ (0,75 điểm)}
\[
z^{(1)}=W^{(1)}x+b^{(1)}=\begin{bmatrix}1\cdot1+0\cdot(-1)\\0\cdot1+(-1)(-1)\end{bmatrix}=\begin{bmatrix}1\\1\end{bmatrix},\qquad
a^{(1)}=\relu(z^{(1)})=\begin{bmatrix}1\\1\end{bmatrix}.
\]
\[
z^{(2)}=W^{(2)}a^{(1)}+b^{(2)}=\begin{bmatrix}1\cdot1+0\cdot1\\0\cdot1+0\cdot1\\0\cdot1+1\cdot1\end{bmatrix}+\begin{bmatrix}0\\0\\0\end{bmatrix}
=\boxed{\begin{bmatrix}1\\0\\1\end{bmatrix}}\in\mathbb R^{3}.
\]

\subsection*{b) Softmax và loss (0,75 điểm)}

Công thức: $\hat y_k=e^{z^{(2)}_k}/\sum_j e^{z^{(2)}_j}$. Với $z^{(2)}=(1,0,1)^{\top}$ và $e=2.718282$:
\[
S=e^1+e^0+e^1=2e+1=6.436564 .
\]
\begin{align*}
\hat y_1&=\frac{e}{S}=\frac{2.718282}{6.436564}=0.422319,\\
\hat y_2&=\frac{1}{S}=\frac{1}{6.436564}=0.155362,\\
\hat y_3&=\frac{e}{S}=0.422319.
\end{align*}
\textbf{Kiểm tra tổng (trước khi làm tròn):}
\[
\hat y_1+\hat y_2+\hat y_3=\frac{e+1+e}{2e+1}=1 \quad(0.422319+0.155362+0.422319=1.000000).
\]
\[
\hat y=\begin{bmatrix}0.4223\\0.1554\\0.4223\end{bmatrix}\ \text{(làm tròn 4 chữ số)}.
\]
Do $y$ là one-hot của lớp 2, chỉ có số hạng $k=2$ khác 0:
\[
L=-\ln\hat y_2=-\ln(0.155362)=\ln(2e+1)=\ln(6.436564)=1.861995\ \Rightarrow\ \boxed{L\approx1.8620}.
\]

\subsection*{c) $\delta^{(2)}$ và gradient tầng ra (1,50 điểm)}

Với softmax kết hợp cross-entropy: $\delta^{(2)}=\pd{L}{z^{(2)}}=\hat y-y$.
\[
\delta^{(2)}=\begin{bmatrix}0.422319\\0.155362\\0.422319\end{bmatrix}-\begin{bmatrix}0\\1\\0\end{bmatrix}
=\boxed{\begin{bmatrix}0.422319\\-0.844638\\0.422319\end{bmatrix}}\approx\begin{bmatrix}0.4223\\-0.8446\\0.4223\end{bmatrix},
\quad\text{kích thước }3\times1.
\]
Vì $z^{(2)}=W^{(2)}a^{(1)}+b^{(2)}$:
\[
\pd{L}{W^{(2)}}=\delta^{(2)}\,(a^{(1)})^{\top}
=\begin{bmatrix}0.422319\\-0.844638\\0.422319\end{bmatrix}\begin{bmatrix}1&1\end{bmatrix}
=\boxed{\begin{bmatrix}0.422319&0.422319\\-0.844638&-0.844638\\0.422319&0.422319\end{bmatrix}}
\]
\[
\approx\begin{bmatrix}0.4223&0.4223\\-0.8446&-0.8446\\0.4223&0.4223\end{bmatrix},
\]
kích thước $3\times2$ (cùng kích thước với $W^{(2)}$).
\[
\pd{L}{b^{(2)}}=\delta^{(2)}=\boxed{\begin{bmatrix}0.4223\\-0.8446\\0.4223\end{bmatrix}},\quad\text{kích thước }3\times1.
\]

\begin{center}
\begin{tabular}{@{}lcc@{}}
\toprule
Đại lượng & Kích thước & Công thức \\
\midrule
$\delta^{(2)}$ & $3\times1$ & $\hat y-y$ \\
$\pd{L}{W^{(2)}}$ & $3\times2$ & $\delta^{(2)}(a^{(1)})^{\top}$ \\
$\pd{L}{b^{(2)}}$ & $3\times1$ & $\delta^{(2)}$ \\
\bottomrule
\end{tabular}
\end{center}

\begin{center}\textbf{--- Hết ---}\end{center}

\end{document}
